\documentclass[letterpaper,10pt,twoside]{amsart}

\usepackage{amsmath,amssymb,amsfonts,amsthm}
%\usepackage{biblatex}
\usepackage{tikz-cd}
\usepackage[margin=0.9in,
left=1in,%
right=1in,%
top=1.25in,%
bottom=1.25in
]{geometry}	
\usepackage{stmaryrd} %For mapsfrom.

\usepackage[colorinlistoftodos]{todonotes} %For todonotes.
\usepackage{tcolorbox}

\usepackage{lscape}

\usepackage{xstring}


%\usepackage[affil-it]{authblk}
%\usepackage{setspace} %For good space between ToC entries.

%\addtolength{\topmargin}{-.875in}

%%%%------ Fonts-------------
%\usepackage{fontspec}
%\setmainfont{QTHandwriting}
% \usepackage{boisik}
% \usepackage[OT1]{fontenc}

% \usepackage{mlmodern}
% \usepackage[T1]{fontenc}
% \usepackage{palatino}
%--------------------------------


%\usepackage{quiver}
\usepackage{bm}
\usepackage{fancyhdr}
\usepackage{mathrsfs}
\usepackage{amsbsy}
%\usepackage{titlesec}
\usepackage{yhmath}
\usepackage{pifont}

%Hyperref Settings------
\usepackage{hyperref}
\usepackage{xcolor}
\hypersetup{
	colorlinks=true,
	linkcolor={blue},
	citecolor={blue},
	anchorcolor={blue}
}

%Macro for acts in proofs
\newenvironment{act}[2]{\begin{center}
		\textbf{Act #1} : \textit{#2}
	\end{center}
}

%Font packages:-


\usepackage[scr=euler,scrscaled=1.05]{mathalfa} %For different symbol for Categories
\usepackage{dashbox}
\usepackage{euscript}

%TITLE:-
\title{Functor of points and algebraic vector bundles}
\author[]{Thomas Arnstein}
\date{\today}

%Header and Footer style:-
%\pagestyle{fancy}
%\fancyhf{}
%\fancyhead[CO]{{\fontfamily{lmss}\selectfont \ssmall \rightmark}}
%\fancyhead[CE]{{\fontfamily{lmss}\selectfont \ssmall \leftmark}}%Add name instead of date.
%\fancyhead[LE,RO]{{\fontfamily{lmss}\selectfont \ssmall \thepage}}
%\fancyhead[LO,RE]{}
%\fancyfoot[CE,CO]{}
%\fancyfoot[LE,RO]{}
\setlength{\headheight}{15pt}
\renewcommand{\headrulewidth}{0pt}

\usepackage{moresize}



%MATH BACKGROUND DECLARATORS-------------------------------------------------------
\theoremstyle{proposition}
\newtheorem{proposition}{Proposition}[subsection]

\theoremstyle{definition}
\newtheorem{definition}[proposition]{Definition}

\theoremstyle{theorem}
\newtheorem{theorem}[proposition]{Theorem}

\theoremstyle{definition}
\newtheorem{remark}[proposition]{\textbf{Remark}}

\theoremstyle{definition}
\newtheorem{notation}[proposition]{\textbf{Notation}}

\theoremstyle{definition}
\newtheorem{discussion}[proposition]{\textbf{Discussion}}

\theoremstyle{lemma}
\newtheorem{lemma}[proposition]{\textbf{Lemma}}

\theoremstyle{definition}
\newtheorem{example}[proposition]{\textbf{Example}}

%\theoremstyle{remark}
%\newtheorem*{comment}{\textbf{Comments on Proof Technique}}

\theoremstyle{definition}
\newtheorem{construct}[proposition]{Construction}

\theoremstyle{corollary}
\newtheorem{corollary}[proposition]{Corollary}

\theoremstyle{definition}
\newtheorem{caution}[proposition]{\textbf{Caution}}

\def\Hom{\text{Hom}}
\def\Spec{\text{Spec}}
\def\O{\mathcal{O}}
\def\Proj{\text{Proj}}
\def\P{\mathbb{P}}
\def\A{\mathbb{A}}
\def\F{\mathcal{F}}

%---------------------------------------------------------------------------------
%\setcounter{tocdepth}{2}




\begin{document}
	\maketitle
    \begin{abstract}
        We recall the functor of points perspective for schemes and show that $Sch$ embeds fully faithfully into the functor category $Fun(CRing,Set)$. We'll look at projective schemes analogously. We recall quasi-coherent sheaves on a scheme and begin to look at a classification of vector bundles on $\mathbb{P}^1$.
    \end{abstract}
	\tableofcontents

\section{Schemes as locally ringed spaces}

We recall the definition of a scheme as a locally ringed space.

\begin{definition}
    A \textit{ringed space} is a pair $(X,\O_X)$ with $X$ a topological space and $\O_X$ a sheaf of rings on $X$. A \textit{locally ringed space} is a ringed space such that for each $x\in X$, the stalk $\O_{X,x}$ is a local ring.
\end{definition}

\begin{definition}
    A \textit{scheme} is a locally ringed space $(X,\O_X)$ which is locally isomorphic to an affine scheme. That is, every $x\in X$ has is an affine open neighborhood $(U,\O_X|_U) \cong (\Spec(R),\O_{\Spec(R)})$ for a ring $R$.
\end{definition}

\begin{remark}
    \begin{enumerate}
        \item The functor $R\mapsto (\Spec(R),\O_{\Spec(R)})$ is fully faithful from the category of commutative rings to the category of ringed spaces. We think of the category $Aff$ of affine schemes as the subcategory spanned by the functor. 
        \item \cite{EH} For any ring $R$ and any scheme $X$, there is a natural bijection \[Mor(X,\Spec(R))\cong\Hom_{CRing}(R,\O_X(X))\]

    \end{enumerate}
\end{remark}

\section{Functor of points}
The functor of points is a way to view the data of a scheme as a certain functor from rings to sets. So a scheme is an organized collection of sets, and constructions in the category of schemes reduces to understanding which functors some from schemes. In fact we embed the category of schemes into a functor category. Since this functor category is in particular a topos/site, we can use descriptive operations as in $Set$. 

\begin{definition}
    The \textit{functor of points} of a scheme $X$ is the representable functor \[Mor(-,X)=:h_X: Sch^{op}\to Set\] We call the set $Mor(Y,X)=h_X(Y)$ the set of $Y$-valued points on $X$.
\end{definition}
\begin{remark}
    Since this construction is natural in $X$, we get a functor $h:Sch\to Fun(Sch^{op},Set)$. Moreover if we restrict to the category of affine schemes $Aff\simeq (CRing)^{op}$, then $h:Aff\to Fun(Aff^{op},Set) \simeq Fun(CRing,Set)$ is the Yoneda embedding. Importantly, we will see that the functor of points of a scheme is determined by its restriction to affine schemes.
\end{remark}

\begin{theorem}\cite{EH}
Let $k$ be a commutative ring. Then the category $Sch_k$ of $k$-schemes embeds into the functor category $Fun(Aff_k^{op},Set)\simeq Fun(k\text{-alg},Set)$ via the functor $X\mapsto Mor_k(-,X)=h_X$.
\end{theorem}
\begin{proof}
    Let $\varphi: h_X\to h_{X'}$ be a natural transformation. The claim is that $\varphi\cong h_f$ for a unique $k$-morphism $f:X\to X'$. Let $\{U_\alpha\}$ be an affine cover of $X$. Applying $\varphi$ to the inclusion $i_\alpha:U_\alpha\to X$, we define $f|_{U_\alpha} = \varphi_{U_\alpha}(i_\alpha)\in Mor_k(U_\alpha,X')$. This is seen from the diagram \[\begin{tikzcd}
	{h_X(X)} & {} & {h_X(U_\alpha)} \\
	{h_{X'}(X)} & {} & {h_{X'}(U_\alpha)}
	\arrow["{id_X\mapsto i_\alpha}", from=1-1, to=1-3]
	\arrow["{\varphi_X}"', from=1-1, to=2-1]
	\arrow["{\varphi_{U_\alpha}}", from=1-3, to=2-3]
	\arrow["{i_\alpha^*}"', from=2-1, to=2-3]
\end{tikzcd}\] Then define $f = \varphi_X(id_X)$ so that $i_a^*(f)=f|_{U_\alpha}$. The compatibility conditions for $\{U_\alpha\}$ to be a cover exactly mean $f$ is well-defined. This can be seen by taking the above square. It is also clear that $f$ is unique. If $f\neq g: X \rightrightarrows X'$ then $f|_{U_\alpha} \neq g|_{U_{\alpha}}$ for some $\alpha$.
\end{proof} 

This theorem amounts to the statement that morphisms of schemes are determined by their restrictions to affine opens. Moreover, it asserts that we can take limits pointwise in $Fun(Aff_k^{op},Set)$. 

The upshot is that we can completely view schemes as functors $CRing\to Set$. We can define open and closed subfunctors to be those represented by inclusions of open and closed subschemes. 

\section{Quasi-coherent sheaves}

We recall quasi-coherent sheaves on schemes. For a scheme $(X,\O_X)$, an $\O_X$-module is a sheaf $\F$ on $X$ such that for every open $U\subset X$, $\F(U)$ is naturally an $\O_X(U)$-module. 
\begin{definition}
    An $\O_X$-module $\F$ is a \textit{quasi-coherent sheaf} on $X$ if either of the following equivalent conditions hold:
    \begin{enumerate}
    \item For every inclusion of affine opens $V\subset U\subset X$, $\F(U)\otimes_{\O_X(U)}\O_X(V)\to \F(V)$ is an isomorphism of $\O_X(V)$-modules.
    \item For every affine open $U\subset X$, $\F|_U \cong \text{coker}(\bigoplus_I\O_X(U)\xrightarrow{\varphi} \bigoplus_J \O_X(U))$ for some $\O_X(U)$-linear map $\varphi$. ($I$ and $J$ are not assumed to be finite.) Note that this is also called being locally presentable.


    \end{enumerate}
\end{definition}

The category of quasi-coherent sheaves on a scheme $X$ is denoted $QCoh(X)$. This is an abelian category.

Moreover, if $X=\Spec(R)$ is affine, then the category $QCoh(X)$ of quasi-coherent sheaves on $X$ is equivalent to the category of $R$-modules.

\begin{remark}
    If we view a scheme $X$ as a functor $Aff^{op}\to Set$, then a quasi-coherent sheaf on $X$ is an assignment to each $x \in X(\Spec(R))$ an $R$-module $\F_x$ subject to some compatibility. Given a quasi-coherent sheaf $\F$ on $X$ (in the sense of the previous definition), then $\F_x$ is the pullback of $\F$ along $x$.
\end{remark}

Quasi-coherent sheaves on a scheme play the role of vector bundles on a scheme. If $X$ is a $k$-scheme for a field $k$, a quasi-coherent sheaf on $X$ is a locally presentable $k$-module, that is, a $k$-vector space and taking an affine cover of we start to see the local trivializations as we're used to in topology.

\section{Projective space}

In this section we compare perspectives on $\mathbb{P}^n$. Projective space can be described as a gluing of affine spaces, as $\Proj(A[x_0,\dots,x_r])$, or by a functor of points.

\begin{construct}
    Fix a base ring $k$ and let $U=\Spec(k[y])$ and $V=\Spec(k[z])$. There are affine schemes $\A^1_k$ which have the multiplicative group scheme $G_m=\Spec(k[x,x^{-1}])$ as a subscheme. Then $\P^1_k$ is the pushout \[\begin{tikzcd}
	{G_m} & {\A^1\cong U} \\
	{\A^1\cong V} & {\P^1}
	\arrow["{(y\mapsto x)^*}", from=1-1, to=1-2]
	\arrow["{(z\mapsto x^{-1})^*}"', from=1-1, to=2-1]
	\arrow[from=1-2, to=2-2]
	\arrow[from=2-1, to=2-2]
\end{tikzcd}\] This is gluing two copies of $\A^1_k$ along the isomorphism $G_m\to G_m$ given by $x\mapsto x^{-1}$. The coordinate ring of $\P^1$ is then $\O_{\P^1}(\P^1) \cong k[t]\otimes_{k[t,t^{-1}]}k[t^{-1}]\cong k$. This construction can be generalized where $\P^r_k$ is defined by gluing $r+1$ copies of $\A^r_k$. Fixing $r$ and letting $S=k[x_0,\dots,x_r]$, we have for each $0\leq i\leq r$, we have a map \[in_i:U_i:=\Spec(k[x_0,\dots,\hat{x_i},\dots x_r])\cong \A^r_k \to \Spec(S)\cong \A^{r+1}_k\] given on rings by $x_k\mapsto x_k/x_i$. The isomorphism on intersections $U_i\cap U_j$ is given by $x_i\mapsto x_i^{-1}$ and $x_j\mapsto x_j^{-1}$. This is also to say that $\P^r_k = (\A^{r+1}_k\setminus\{0\})/G_m$.
\end{construct}

\begin{construct}
    To any ($\mathbb{Z}$-)graded ring $S = \bigoplus_i S_i$ we associate a set $\Proj(S) = \{p \in \Spec^h(S) : p\not\subset S_+\}$ where $S_+=\bigoplus_{i\geq 1}S_i$ and $\Spec^h$ indicates that $p$ is a homogeneous prime ideal. A homogeneous prime ideal is a prime ideal which is a graded submodule of $S$. The topology on $\Proj(S)$ is generated by the open sets $D^+(f)=\{p\in \Proj(S):f\notin p\}$. Closed subsets are then of the form $V_+(I) = \{p\in\Proj(S): p\supset I\}$ for an homogeneous ideal $I$. This defines a projective scheme for any graded ring $S$. If we have a list $x_0,\dots,x_r \in S$ which generate $S_+$, then the opens $(\Proj(S))_i=\Proj(S)\setminus V(x_i)$ form an affine open cover since $(\Proj(S))_i \cong \Spec(S[x_i^{-1}]_0)$. In particular, when $S=k[x_0,\dots,x_r]$ we have $\Proj(S)=\P^r_k$ and $(\P^r_k)_i\cong \A^r_k$ for all $i=0,\dots,r$.
\end{construct}
\begin{construct}
    Lastly, we describe the functor of points associated to $\P^r_k$. Recall in classical topology, when $k$ is $\mathbb{R}$ or $\mathbb{C}$, then $\P^r_k=\{\text{lines through origin in }k^{r+1}\} = (k^{n+1}\setminus 0)/(\lambda v\sim v)$ for $\lambda \in k^\times$. For a ring $R$, we want that $\P^r$ represents the functor \[Aff^{op}\to Set, \hspace{0.1in} \Spec(R)\mapsto\{L \subset R^{n+1} \text{ projective submodules of rank }1\}\] And for $\Spec(S)\xrightarrow{f}\Spec(R)$ ie $f_\sharp:R\to S$, we take $L\mapsto L\otimes_R S$.
\end{construct}
We now describe morphisms to $\P^r$.

\begin{theorem}
    Let $T$ be a local $k$-algebra. Then the set $Mor_k(\Spec(T),\P^r_k)$ is in bijection with the set of $(r+1)$-tuples $[a_0:\dots:a_r] \in k^{r+1}$ such that $a_i\in k^\times$ for some $0\leq i\leq r$, modulo the equivalence relation \[[a_0:\dots:a_r]\sim[\lambda a_0:\dots:\lambda a_r], \hspace{0.1in} \lambda \in k^\times\]
\end{theorem}

Recall that a sheaf $P$ on $X$ is invertible if the functor \[-\otimes_{\O_X} P: \O_X\text{-mod}\to \O_X\text{-mod}\] is an equivalence. Equivalently, $P$ is invertible if there exists a sheaf $\mathcal{N}$ such that $\mathcal{N}\otimes_{\O_X}P \cong \O_X$. This can be checked locally on affines $U=\Spec(A)$ with $P|_U$ a projective $A$-module and then $\mathcal{N}|_U=\underline{\Hom}_{\O_U}(P|_U,\O_U) = \Hom_A(P|_U,A)$. \cite{Stacks1}

\begin{theorem}\cite{EH}
    For any scheme $X$, we have natural bijections \[Mor(X,\P^r)\cong\{\text{subsheaves }K\subset \O_X^{r+1} \text{ which are locally free of rank }r\}\]
    \[\cong\{\text{invertible sheaves }P \text{ together with epimorphisms }\O_X^{r+1}\to P\}/O_X(X)^\times\]
\end{theorem}
Due to the section 2, we can assume that $X=\Spec(A)$ is affine. An invertible sheaf on $\Spec(A)$ is exactly a projective $A$-module of rank 1. Before the proof, we recall some of this.

\section{Vector bundles on $\mathbb{P}^1$}

Following chapter 8 of \cite{1}.

We fix a ring $k$ and view $\P^1_k$ as the scheme obtained by gluing $\Spec k[t]$ to $\Spec k[t^{-1}]$ along $\Spec k[t,t^{-1}]$. Let $P_+$ be a projective $k[t]$-module and $P_-$ a projective $k[t^{-1}]$-module, both of rank $r$. These are quasi-coherent sheaves (vector bundles) on standard the affine cover of $\P^1_k$. By Quillen-Suslin, $P_+$ and $P_-$ are actually free. By restricting along the structure maps, we obtain two projective $k[t,t^{-1}]$-modules \[\widetilde{P_+}= P_+\otimes_{k[t]}k[t,t^{-1}] \text{ and } \widetilde{P_-}=P_-\otimes_{k[t^{-1}]}k[t,t^{-1}]\] The compatibility of the structure maps actually gives an isomorphism $\widetilde{P_+}\xrightarrow{\varphi} \widetilde{P_-}$. If we let $\{e_i^+\}$ and $\{e_i^-\}$ be bases of $P_+$ bzw $P_-$, then the choice of an isomorphism $\varphi$ is determined by a linear transformation $A \in GL_r(k[t,t^{-1}])$. This is determined up to choice of basis. Base change of $P_+$ corresponds to a left action by $GL_r(k[t])$ and base change of $P_-$ a right action by $GL_r(k[t^{-1}])$. This defines an injective function \[GL_r(k[t])\backslash(GL_r(k[t,t^{-1}])/GL_r(k[t^{-1}]) \to Vect_r(\P^1_k)\] where $Vect_r(\P^1_k)$ is the set of isomorphism classes of rank $r$ vector bundles on $\P^1_r$.

\begin{theorem}\cite{1}
    If $k$ is a field, then the above map is a bijection \[GL_r(k[t])\backslash(GL_r(k[t,t^{-1}])/GL_r(k[t^{-1}]) \cong Vect_r(\P^1_k)\] A matrix representation of a vector bundle $\F$ on $\P^1_k$ is a clutching function.
\end{theorem}

\begin{remark}
    In particular, the theorem states that isomorphism classes of line bundles over $\P^1_k$ are in bijection to $k[t]^\times\backslash k[t,t^{-1}]^\times /k[t^{-1}]^\times$. Every unit in $k[t,t^{-1}]$ can be written as $at^r$ for $a\in k^\times$ and $r\in\mathbb{Z}$. The line bundle on $\P^1_k$ represented by $(t^{-d})$ is written as $\O(d)$. As a sheaf, the sections $\Gamma(\P^1_k,\O(d))\cong k[x_0,x_1]_d$ the $d$th graded piece of the graded ring for which $\P^1_k=\Proj(k[x_0,x_1])$. 
\end{remark}

\begin{theorem}\cite{1}
    If $k$ is a field, then every vector bundle $\F$ on $\P^1_k$ can be written as a direct sum of line bundles \[\F\cong \O(a_1)\oplus\dots\oplus \O(a_n)\] for some integers $a_1\geq\dots\geq a_n$. The clutching function of $\F$ is given by $\text{diag}(t^{a_1},\dots,t^{a_n})$.
\end{theorem}

Let $R=k[x]$ and consider the scheme $\P^1_k\times_k \A^1_k = \Proj(k[x_0,x_1])\times_{\Spec(k)}\Spec(R) = \P^1_k(R) = \P^1_R$. We want to show here the failure of $\A^1$-invariance of vector bundles. That is, we show that $Vect_2(\P^1_k)\not\cong Vect_2(\P^1_k\times_k \A^1_k)$, despite $\A^1$ having no nontrivial vector bundles.

Let $x$ be the coordinate for $\A^1$ and $t$ the coordinate for $\P^1$. Consider the rank 2 vector bundle $\F$ on $\P^1(R)$ given by the clutching function \[\begin{pmatrix} t & 0 \\ x & t^{-1} \end{pmatrix}\] Let $\pi:\P^1\times_k\A^1\to \P^1$ denote the canonical map from the pullback, and we consider the map \[\pi^*:Vect_2(\P^1)\to Vect_2(\P^1\times_k \A^1)\] given by pulling back a vector bundle along $\pi$. We show that this bundle is not obtained by pulling back a bundle on $\P^1_k$. When $x=0$, this vector bundle is clearly $\O(1)\oplus \O(-1)$. At $x=1$, we have \[\begin{pmatrix} 0 & 1 \\ -1 & t \end{pmatrix}\begin{pmatrix} t & 0 \\ 1 & t^{-1} \end{pmatrix}\begin{pmatrix} 1 & -t^{-1} \\ 0 & 1 \end{pmatrix} = \begin{pmatrix} 1 & 0 \\ 0 & 1 \end{pmatrix}\] so that the fiber over $x=1$ is the vector bundle $\O(0)\oplus \O(0)$. We therefore conclude that $\F$ is not pulled back from a vector bundle on $\P^1$, so the map $\pi^*$ is not surjective. 

This illustrates our first example of the failure of $\A^1$-invariance of vector bundles.

\bibliography{ref}{}
\bibliographystyle{alpha}

\end{document}